\section{A surface relation from determinant-fixed Quot schemes}
\label{sec:quot}

The tensor formalism of Section~\ref{sec:tensors} does not by itself bound
products of the odd coefficients of a projected small diagonal. We obtain
the additional relation we need from the surface square of the curve.
The decisive features are a determinant constraint, which removes the
surface obstruction directions, and a marked Hilbert-scheme calculation,
which retains the Chow classes of all output positions.

\subsection{The surface bound}

Let \(\pi_C:C\to S\) be a smooth projective family of connected curves of
genus \(b\), with a section, over an irreducible smooth quasi-projective
complex variety. Choose a rational divisor \(e_C\) of relative degree one,
and use the corrected divisor \(e'_C\), projector \(\eta\), and odd
alternating cubic form \(\varphi\) of Proposition~\ref{prop:tensor-test}.
Set
\begin{equation}
\label{eq:surface-parameters}
 \chi=(b-1)^2,\qquad s_K=8\chi,\qquad d=3,\qquad
 t_0=\chi+\frac{d(d-1)}2s_K-1=25\chi-1,\qquad m=3t_0.
\end{equation}
Throughout this section, the support quotient is the quotient by Chow
classes that vanish after deleting a closed subset of \(S\) of codimension
at least \(b-1\). Its coefficient superalgebra will be denoted by
\(\mathfrak B=\mathfrak B_{S,b-1}\).

\begin{theorem}[Surface bound]
\label{thm:surface}
Suppose \(b\geq5\) and \(\dim S<\chi-1\).
Let \((w_\alpha)_{\alpha\in I}\) be any finite collection of vectors in the
odd symplectic standard space \(W_b\) whose span is isotropic. For any
matrix of scalars \(\gamma=(\gamma_{\alpha\beta})_{\alpha,\beta\in I}\),
put
\begin{equation}
\label{eq:surface-tensor}
 T_\gamma=
 \sum_{\substack{\alpha_1,\alpha_2,\alpha_3\in I\\
                  \beta_1,\beta_2,\beta_3\in I}}
 \left(\prod_{j=1}^3\gamma_{\alpha_j\beta_j}\right)
 \varphi(w_{\alpha_1},w_{\alpha_2},w_{\alpha_3})
 \varphi(w_{\beta_1},w_{\beta_2},w_{\beta_3}).
\end{equation}
Then
\begin{equation}
\label{eq:surface-bound}
 T_\gamma^{\,t_0}=0\qquad\text{in }\mathfrak B.
\end{equation}
The assertion remains valid after any characteristic-zero scalar
extension of the tensor formalism.
\end{theorem}

Isotropic tests alone do not separate arbitrary Chow tensors.
In Lemma~\ref{lem:large-vertex-test}, we first fix an arbitrary color
map, find a subspace of its domain with isotropic image, and use the
exterior-ideal argument to annihilate the entire sufficiently long
product, throughout modulo base support of codimension at least
\(b-1\). Only after doing this for every color map do we apply the
faithful test of Proposition~\ref{prop:tensor-test}.

We prove the localization relation underlying this theorem in the present
section. Section~\ref{sec:coefficient} proves that its scalar coefficient
is nonzero and thereby completes the proof.

Write
\[
 X=C\times_S C,\qquad \pi_X:X\longrightarrow S,\qquad
 K=K_{X/S}.
\]
We write line bundles additively. On every fiber, \(K\) is ample,
\(K^2=s_K\), and \(\chi(\mathcal O_X)=\chi\).
The curve section gives a section \(o:S\to X\).
The two summands of the target bundle and their torus weights are
\begin{equation}
\label{eq:quot-target}
 V=L_1\oplus L_2=K\oplus\mathcal O_X,\qquad
 (w_1,w_2)=(0,1),\qquad A_0=L_1-L_2=K.
\end{equation}
Here the use of \(w_1,w_2\) for weights is confined to formulas concerning
the target bundle; the tensor vectors retain their indices \(\alpha\).
Let \(\mathcal Q\) be the relative Quot scheme of rank-zero quotients
\[
 0\longrightarrow I'\longrightarrow V\longrightarrow Q
 \longrightarrow0
\]
with the following determinant and numerical data:
\begin{equation}
\label{eq:quot-data}
 D=dK,\qquad
 \det Q=D\text{ in }\Pic(X/S),\qquad
 \int_{X_s}\ch_2(Q_s)=\nu+L_1D-\frac{D^2}{2},
 \qquad \nu=t_0-1.
\end{equation}
All intersections in this display are fiber intersections.
One may use the relative polarization \(K\) to specify the Hilbert
polynomial. Projective representability and the Grassmannian construction
are supplied by Grothendieck's Quot theorem
\cite[Theorem~5.1 and Section~5]{Nitsure}. The determinant condition is a closed fiber of the separated relative
Picard scheme \cite[Theorem~4.8(1)]{KleimanPicard}, so \(\mathcal Q\)
is projective over \(S\).

\subsection{The derived obstruction theories}

The determinant constraint changes both the obstruction space and the
virtual dimension of \(\mathcal Q\).  We first construct its tangent
complex as the fiber of the determinant differential.  The surjectivity
of that differential on obstruction spaces will give a two-term theory.
The same construction on the torus-fixed factors will then express their
virtual classes as virtual pullbacks from one rank-one Quot scheme.

\subsubsection{Arrows and the Quot tangent}

The perfect-complex moduli theorem of
\cite[Corollaries~3.15, 3.17, and~3.19, Definition~3.18,
Lemma~3.27, and Corollary~3.29]{TV}
applies over every affine open of \(S\).
Indeed, the dg-category of perfect complexes on a smooth proper scheme
over an affine base is proper by perfection of relative derived Hom
and smooth by perfection of its diagonal. These smooth proper categories are saturated in the sense of
\cite[Lemma~3.27 and Corollary~3.29]{TV}. Their moduli of perfect objects
and perfect arrows are locally geometric and locally of finite
presentation. These constructions glue over \(S\).
The tangent of the object moduli at a perfect complex \(F\) is
\(\mathbf R\Hom(F,F)[1]\).
All Hom complexes in what follows are derived Hom on the surface,
followed by derived pushforward to the parameter scheme. Perfect direct
images and derived base change are used in the form of
\cite[Tag~0B91]{Stacks}, and relative duality in the form of
\cite[Tags~0BRT, 0A9Q, and~0B6I]{Stacks}.

\begin{lemma}[The open arrow enhancement]
\label{lem:quot-arrow}
The derived stack of arrows with target \(V\) has an open derived
substack whose classical truncation is the ordinary relative Quot scheme
with the prescribed Hilbert polynomial, before fixing determinant.
On this truncation its relative tangent complex is
\[
 \operatorname{Cone}\bigl(
 \mathbf R\Hom(I',I')\longrightarrow
 \mathbf R\Hom(I',V)\bigr)
 \simeq \mathbf R\Hom(I',Q).
\]
This is a perfect complex of Tor-amplitude \([0,1]\), compatible with
base change. The differential of forgetting the arrow and keeping its
source is the boundary map
\[
 \mathbf R\Hom(I',Q)\longrightarrow\mathbf R\Hom(I',I')[1].
\]
\end{lemma}

\begin{proof}
In the perfect-arrow stack take the locus where the source and cone are
degree-zero sheaves on surface fibers, with the prescribed Hilbert data.
For a perfect complex, relative Tor-amplitude zero is open.
Concretely, a locally free resolution can be truncated first at its
positive end; at its negative end, fiberwise injectivity and the local
flatness criterion successively give flat cokernels.
Properness of the surface family makes the condition open on the
parameter scheme. Thus on a classical base the selected arrows are
injections with flat sheaf quotients.

Conversely, the kernel and quotient of a flat quotient family are perfect
on the total surface family, including over a singular parameter scheme.
In a local free resolution the successive syzygies remain flat over the
parameters. Regularity of the surface fibers makes the second syzygy
fiberwise locally free; the local flatness criterion then makes it
locally free on the total space. Hence the resolution terminates.
These arguments commute with base change. An automorphism of the source
over its inclusion in \(V\) is the identity, and negative Ext groups of
the sheaves vanish. The classical moduli problem is therefore the Quot
scheme, not a gerbe over it.

For an arrow \(F\to G\) with both vertices variable, the endomorphism
complex of the diagram is
\[
 \operatorname{fib}\bigl(
 \mathbf R\Hom(F,F)\oplus\mathbf R\Hom(G,G)
 \longrightarrow\mathbf R\Hom(F,G)\bigr),
\]
where the map is the difference of pre- and postcomposition.
Its shift by \([1]\) is the tangent of the arrow moduli.
Taking the tangent fiber over the fixed target gives the displayed
cone and its boundary map. These identifications are identifications of
perfect relative complexes: they can be computed with a locally free
resolution of the one-arrow diagram and commute with derived base
change.

At a geometric point, surface Serre duality gives
\[
 \Ext^2(I',Q)^\vee=\Hom(Q,I'(K))=0,
\]
since \(Q\) is torsion and \(I'\) is torsion-free.
The other possible Ext degrees lie in \([0,1]\).
The geometric-fiber criterion for a perfect complex proves the asserted
Tor-amplitude.
\end{proof}

\subsubsection{Fixing the relative determinant}

Let \(\mathcal P_o^{\mathrm{der}}\) be the derived line stack
rigidified along \(o\), namely the fiber of evaluation along \(o\) over
the trivial line. For any line \(M\) on a base-changed surface, its
normalization is
\[
 M\otimes\pi_X^*(o^*M)^{-1}.
\]
It has its canonical trivialization along \(o\).
Because \((\pi_X)_*\mathcal O_X=\mathcal O_S\) universally, all
line automorphisms come from the base, and the normalization removes
them. The classical truncation of \(\mathcal P_o^{\mathrm{der}}\) is
therefore the relative Picard scheme, parametrizing lines modulo base
twists. Its tangent is
\begin{equation}
\label{eq:based-picard-tangent}
 \mathbf R\pi_{X*}\mathcal I_o[1],
\end{equation}
where constant functions split evaluation
\(\mathbf R\pi_{X*}\mathcal O_X\to\mathcal O_S\).
This complex has geometric-fiber cohomology \(H^1(\mathcal O_X)\)
in degree zero and \(H^2(\mathcal O_X)\) in degree one.

The determinant of perfect complexes is a functorial derived morphism
\cite[Section 3.1]{STV}. Its differential is trace
\cite[Proposition 3.2]{STV}. After normalization it maps the open
arrow enhancement to \(\mathcal P_o^{\mathrm{der}}\).
Take its derived fiber at the normalized line \(L_1+L_2-D\).
Its truncation is exactly \(\mathcal Q\): it fixes the Picard class
of \(\det I'\), not a trivialization of the unnormalized determinant
on a parameter scheme. The relative tangent restricted to \(\mathcal Q\)
is
\begin{equation}
\label{eq:quot-fixed-tangent}
 T_{\mathcal Q/S}^{\mathrm{vir}}=
 \operatorname{fib}\left\{
 \mathbf R\Hom(I',Q)
 \xrightarrow{\ \partial,\ \mathrm{tr},\ \mathrm{based}\ }
 \mathbf R\pi_{X*}\mathcal I_o[1]\right\}.
\end{equation}
Here the trace identification is an identity of relative complexes,
not an inference from equality at closed points.
Locally represent a perfect family by a bounded complex \(P^\bullet\)
of finite free modules. Its determinant is
\(\bigotimes_j(\det P^j)^{(-1)^j}\).
Differentiation on infinitesimal automorphisms gives the supertrace
\(\sum_j(-1)^j\mathrm{tr}_{P^j}\).
The same calculation with square-zero complex coefficients gives
the chain map \(\mathbf R\Hom(P^\bullet,P^\bullet)\to\mathcal O\);
it respects homotopies and change of perfect resolution.
Using the identification of object tangents with shifted loop
tangents gives its shift by \([1]\) as the determinant differential.
This construction glues and commutes with base change, and therefore
gives the relative trace in \eqref{eq:quot-fixed-tangent}, including
over nonreduced parameter schemes.

\begin{lemma}[The determinant obstruction map]
\label{lem:det-surjective}
Let \(0\to F\to V_0\to Q_0\to0\) be an exact sequence on a smooth
projective surface, where \(V_0\) is a vector bundle of rank \(r\)
and \(Q_0\) is torsion.
If \(K\) is ample and \(Kc_1(Q_0)>rK^2\), the map
\[
 \Ext^1(F,Q_0)\longrightarrow\Ext^2(F,F)
 \xrightarrow{\mathrm{tr}}H^2(\mathcal O_X)
\]
is surjective.
\end{lemma}

\begin{proof}
By duality, a canonical form \(\omega\) annihilating the image has the
property that \(\omega\id_F:F\to F(K)\) extends to a map
\(a:V_0\to F(K)\).
If \(i:F\to V_0\), the maps \(i(K)a\) and
\(\omega\id_{V_0}\) agree where \(F=V_0\), hence everywhere.
It follows that \(\omega Q_0=0\).
For nonzero \(\omega\), write \((\omega)=\sum a_jC_j\).
At the generic point of \(C_j\) a quotient of
\((\mathcal O/(\omega))^r\) has length at most \(ra_j\).
These generic lengths are the multiplicities of \(c_1(Q_0)\).
Thus \(c_1(Q_0)\leq r(\omega)\) as divisorial cycles, and ampleness
gives \(Kc_1(Q_0)\leq rK^2\), a contradiction.
The argument permits embedded points and nonreduced curve support.
\end{proof}

Applying Lemma~\ref{lem:det-surjective} with \(r=2\) and \(D=3K\)
shows that the only possible degree-two cohomology of
\eqref{eq:quot-fixed-tangent} vanishes. The complex is perfect of
Tor-amplitude \([0,1]\).
The canonical map from the cotangent of a derived enhancement to that
of its classical truncation has cone in degrees at most \(-2\)
\cite[Proposition 1.2 and Corollary 1.4]{STV}.
Consequently the dual of \eqref{eq:quot-fixed-tangent} is a relative
perfect obstruction theory in the sense of \cite[Definitions 4.4 and
5.1]{BF}. Adding the tangent of the smooth base by transitivity gives
the absolute theory. All constructions are torus equivariant.

\subsubsection{Virtual rank and compatible pullbacks}

Riemann--Roch gives the relative virtual rank
\begin{equation}
\label{eq:quot-rank}
 v=\chi(I',Q)+\chi-1
   =2\nu+(A_0-K)D+\chi-1
   =2t_0+\chi-3>2t_0.
\end{equation}
Indeed, if \(q_2=\int\ch_2Q\), then
\(\chi(V,Q)=2q_2-(K+L_1+L_2)D\) and
\(\chi(Q,Q)=-D^2\), while the rank of
\eqref{eq:based-picard-tangent} is \(1-\chi\).
Thus fixing determinant contributes the extra \(\chi-1\) in
\eqref{eq:quot-rank}.  The inequality \(v>2t_0\) will force the
marked localization relation.

To compare the fixed virtual class with simpler classes, we need
morphisms of obstruction theories, rather than only tangent classes
in K-theory.  The derived enhancements supply these morphisms as follows.
For a morphism of these derived enhancements, transitivity maps to the
transitivity triangle on truncations. If the two absolute comparison
cones are \(C_Y,C_Z\), the relative comparison cone fits into
\[
 Lf^*C_Z\longrightarrow C_Y\longrightarrow C_{Y/Z}
 \longrightarrow Lf^*C_Z[1].
\]
The absolute cones lie in degrees at most \(-2\); derived pullback
preserves this upper bound and \([1]\) lowers it.
Therefore \(C_{Y/Z}\) also lies in degrees at most \(-2\).
Whenever the relative derived cotangent has perfect amplitude
\([-1,0]\), this diagram supplies the compatible triple of obstruction
theories. Virtual-pullback functoriality then applies
\cite[Definition 4.5, Theorem 4.8, and Corollary 4.9]{Manolache};
see also \cite[Proposition 1.5]{STV}.

\subsection{Fixed loci and their virtual classes}

The equivariant embedding and resolution hypotheses of virtual
localization \cite[Section 1]{GP} hold here.
The Grassmannian construction of relative Quot is equivariant and
embeds the determinant fiber in a smooth scheme over the smooth base
\(S\). One can construct a global equivariant two-term presentation of
its perfect obstruction theory using a linearized ample line.
For completeness, over an affine base sufficiently positive twists
generate its degree-zero cohomology and kill the higher cohomology
of its degree-minus-one cohomology. Lifting generators gives a bundle
mapping to the complex and surjective on degree-zero cohomology;
its derived fiber is a bundle by the Tor-amplitude criterion.
Further Serre vanishing lifts the map to the two-term embedding
cotangent presentation. Torus linear reductivity makes the choices
equivariant.
For a nonaffine \(S\), the same argument may be made after a
Jouanolou affine-bundle pullback
\cite[Proposition~4.1.4(1)]{ADOJouanolou}. The resulting identity descends by
equivariant Chow homotopy invariance and smooth pullback of the
base-changed obstruction theories. Subsequent dimension arguments
are on the original base \(S\).

A torus-fixed subsheaf of \(L_1\oplus L_2\) splits into its two weight
subsheaves, also in families. Kernels and quotients in the two weights
are flat, since they are direct summands of flat families.
Thus a fixed point is a pair of rank-one arrows, subject to the total
determinant constraint. Their direct-sum morphism of derived
enhancements identifies its tangent triangle with the weight-zero part
of \eqref{eq:quot-fixed-tangent}, compatibly with the cotangent
comparison. Normalizing a determinant cancels all character lines
pulled from the base, so the based condition does not acquire a moving
term. The moving virtual normal complex is
\begin{equation}
\label{eq:moving-normal}
 \bigoplus_{i\ne j}\mathbf R\Hom(I'_i,Q_j)
\end{equation}
with weight \(w_j-w_i\).

We use the divisor-and-residual-point factorization established in the
relative setting by \cite[Theorem~1.4]{Fogarty}. We spell out its
compatibility with the determinant and nonreduced bases.
Each rank-one kernel has a unique factorization
\begin{equation}
\label{eq:fixed-factorization}
 I'_i=L_i(-\mathcal C_i)\mathcal I_{Z_i},
 \qquad D_i=[\mathcal C_i]\text{ in }\Pic(X/S),
 \qquad D_1+D_2=D.
\end{equation}
Here \(\mathcal C_i\) is a relative effective Cartier divisor and \(Z_i\)
is a finite-flat point family.
This holds over nonreduced parameter schemes. To see this, let \(T\)
be such a parameter scheme and untwist the target line.
The \(T\)-flat kernel ideal \(I\) has a length-one locally free
resolution, by the fiberwise projective-dimension bound.
Its determinant is a line bundle on \(X_T\), commuting with arbitrary
base change. The maximal-minor map \(I\to\det I\) restricts on every
surface fiber to the inclusion of a rank-one torsion-free sheaf into
its double dual. It is therefore fiberwise injective with
zero-dimensional cokernel. Since both its source and target are
\(T\)-flat, the local flatness criterion makes the map injective
and its cokernel \(\mathcal T\) \(T\)-flat. Properness and the fiber
support make \(\mathcal T\) finite over \(T\).

The inclusion into the original target line extends uniquely across
that finite support: on a smooth relative surface a finite-flat
module has grade at least two, as follows by lifting a fiber regular
sequence with flat successive quotients. Thus its Hom and first Ext
into the structure sheaf vanish. Applying Hom to
\(0\to I\to\det I\to\mathcal T\to0\) gives the extension.
The extension is fiberwise an injection of lines with Cartier
quotient; the local flatness criterion makes its divisor relative
Cartier. It gives the required point ideal and divisor, uniquely and
compatibly with base change.

Let the point lengths be \(r,s\). The Chern character of the two
quotients gives
\begin{equation}
\label{eq:fixed-lengths}
 r+s=\nu+(A_0-D_1)D_2.
\end{equation}
The fixed scheme has only finitely many numerical choices.
Its summand Hilbert polynomials are locally constant.
In addition, the determinant line of each perfect kernel is a
base-change-compatible line family, so its fiber Hilbert polynomial
is locally constant. The residual point length is the difference
between the Euler characteristics of this determinant and its
kernel, and is therefore locally constant as well.
Finite type now gives finiteness of these choices.
All divisor and Picard moduli
used in a choice may be restricted to the required finite-type
quasi-projective pieces. For instance, lines of a fixed Hilbert
polynomial form an open in the moduli of stable rank-one sheaves;
rigidification at \(o\) represents their Picard classes.

Choose a \emph{high} index with \(KD_i>K^2\), which exists since
\(KD_1+KD_2=3K^2\), and call the other index \emph{low}.
Write \(P=D_{\mathrm{low}}\).  For this numerical choice, denote the fixed locus by \(\mathcal F\)
and the unrestricted low rank-one Quot scheme by
\(\mathcal Q_{\mathrm{low}}\).  Forgetting the high arrow gives a map
\[
 f:\mathcal F\longrightarrow\mathcal Q_{\mathrm{low}}.
\]
Its relative derived tangent is
\[
 \operatorname{fib}\left\{
 \mathbf R\Hom(I'_{\mathrm{high}},Q_{\mathrm{high}})
 \longrightarrow\mathbf R\pi_{X*}\mathcal I_o[1]\right\}.
\]
Lemma~\ref{lem:det-surjective}, now with rank one, gives amplitude
\([0,1]\). The transitivity comparison above is therefore a compatible
triple.  Consequently
\begin{equation}
\label{eq:fixed-low-pullback}
 [\mathcal F]^{\mathrm{vir}}
   =f^![\mathcal Q_{\mathrm{low}}]^{\mathrm{vir}}.
\end{equation}
Here \(f^!\) is a bivariant virtual pullback.  We will use its action
on cycle representatives of the low virtual class, including cycles
supported on special divisor loci.

Obstruction Euler factors of point Hilbert schemes also enter the
surface Quot calculations of Oprea and Pandharipande
\cite[Lemma~34 and Section~5.2]{OPQuot}. We establish the relative
comparison with varying determinant and arbitrary irregularity needed
here.

\begin{lemma}[Point excess]
\label{lem:point-excess}
The unrestricted low rank-one Quot scheme is classically
\[
 \operatorname{Div}_{P}(X/S)\times_S X^{[r_{\mathrm{low}}]}_{/S},
\]
where the divisor scheme permits its line to vary in the prescribed
Picard data. Its virtual class is the divisor virtual class times the
smooth relative point-Hilbert class, capped with
\begin{equation}
\label{eq:point-obstruction}
 e(F_{\mathrm{low}}),\qquad
 F_i=\left(\pi_*\mathcal O_{Z_i}(K-\mathcal C_i)\right)^\vee.
\end{equation}
\end{lemma}

\begin{proof}
The divisor enhancement parametrizes invertible-source arrows.
Its tangent is
\(\mathbf R\pi_*\mathcal O_{\mathcal C}(\mathcal C)\), of amplitude
\([0,1]\).
Tensoring its universal source with the universal point ideal gives
a morphism from this enhancement times the smooth relative Hilbert
scheme to the rank-one arrow enhancement, with identity truncation.
The smoothness of the relative point Hilbert scheme follows from the
smooth surface Hilbert theorem \cite[Theorem 2.4]{Fogarty}.
Indeed, choose disjoint etale surface charts around the finite support,
with etale maps to the relative affine plane. Finite-subscheme
deformations lift uniquely through these etale maps. Thus locally on
the Hilbert functor the question is the product of the corresponding
plane Hilbert functors over \(S\), which is smooth over \(S\).

Compare first the source objects of the arrows. The scalar trace
summand of \(\mathbf R\Hom(\mathcal I_Z,\mathcal I_Z)\) splits off.
Writing \(A=\mathbf R\pi_*\mathcal O_X\) and denoting the trace-free
summand by \(B_Z\), the source-object tangent is therefore
\(A[1]\oplus B_Z[1]\). The following calculation identifies
\(B_Z[1]\) with the relative Hilbert tangent, while retaining the
entire Picard summand \(A[1]\).
The trace-free Ext groups in degrees zero and two vanish: an
endomorphism extends to the determinant line, and surface duality
gives the degree-two assertion. Riemann--Roch leaves a trace-free
degree-one space of dimension \(2r_{\mathrm{low}}\).
The point-ideal boundary injects
\(\Hom(\mathcal I_Z,\mathcal O_Z)\) into this space, since every
map \(\mathcal I_Z\to\mathcal O_X\) is a scalar inclusion with zero
composite into \(\mathcal O_Z\).
This boundary is the Hilbert Kodaira--Spencer map; its trace is zero
because the point ideal has trivial determinant. Smoothness and the
dimension \(2r_{\mathrm{low}}\) make it an isomorphism.
The family complexes are perfect, so the fiberwise isomorphism is an
isomorphism of complexes after including the scalar line-moduli
tangent.

For the arrow directions, the comparison is
\[
 \mathbf R\Hom(\mathcal O_X,\mathcal O_X(\mathcal C))
 \longrightarrow
 \mathbf R\Hom(\mathcal I_Z,\mathcal O_X(\mathcal C)).
\]
Its cone is
\[
 \mathbf R\Hom(\mathcal O_Z,\mathcal O_X(\mathcal C))[1]
 \simeq F_{\mathrm{low}}[-1]
\]
by relative Serre duality. The latter is a genuine vector bundle in
degree one: \(Z\) is finite flat and its direct image has no higher
cohomology.
We have therefore obtained a morphism of tangent theories whose cone
is this excess obstruction bundle.

Because it comes from the morphism of enhancements, the dual
comparison commutes with both obstruction morphisms to the same
classical cotangent complex.
The smaller obstruction bundle stack embeds in the larger with
quotient \(F_{\mathrm{low}}\), and the intrinsic-cone embedding
factors through the smaller stack. Zero-section excess intersection
therefore multiplies the smaller virtual class by
\(e(F_{\mathrm{low}})\). This proves the assertion, including when
the divisor scheme is singular.
\end{proof}

Let
\[
 \rho:\mathcal Q_{\mathrm{low}}
   =\operatorname{Div}_{P}(X/S)\times_S X^{[r_{\mathrm{low}}]}_{/S}
   \longrightarrow\operatorname{Div}_{P}(X/S)
\]
be the smooth projection.  Combining point excess with
\eqref{eq:fixed-low-pullback} gives the fixed-cycle formula
\begin{equation}
\label{eq:fixed-cycle-formula}
 [\mathcal F]^{\mathrm{vir}}
 =f^!\left(
   e(F_{\mathrm{low}})\cap
   \rho^*[\operatorname{Div}_{P}(X/S)]^{\mathrm{vir}}
       \right).
\end{equation}
This formula separates the two remaining tasks.  We first determine
where the low divisor virtual cycle can be supported.  Over the
surviving support, the high map \(f\) becomes smooth, and its virtual
pullback will supply the second point Euler factor.

\subsection{Support of the low divisor class}

Put \(y=KP/s_K\).
The virtual relative dimension of the divisor scheme is
\begin{equation}
\label{eq:divisor-virtual-dimension}
 \frac{P(P-K)}2.
\end{equation}
If \(0<y<1\), this is at most \(-(b-1)\).
To check the integer bound, put \(q=b-1\).
The positive intersection \(p=KP\) is a multiple of \(2q\):
each of the two canonical divisors on the product has this fiber
degree. Thus
\[
 2q\leq p\leq8q^2-2q,\qquad
 \frac{P(P-K)}2
 \leq\frac{p^2/(8q^2)-p}{2}
 \leq-q+\frac14
\]
by Hodge index and convexity in \(p\).
The left side is integral, proving the claim.
The absolute divisor virtual cycle consequently has a representative
of dimension at most \(\dim S-(b-1)\).  The image of its support in
\(S\) has codimension at least \(b-1\).  In
\eqref{eq:fixed-cycle-formula}, smooth pullback by \(\rho\) and cap
product with \(e(F_{\mathrm{low}})\) preserve this base support.
So does \(f^!\): it is bivariant and commutes with restriction and
proper pushforward \cite[Theorems 4.1 and 4.4]{Manolache}.
Although these pullbacks can increase the dimension of a cycle, they
do not enlarge the locus over which it lies in \(S\).  Every marked
fixed contribution above this low cycle therefore vanishes in our
support quotient.

On the Picard open where \(H^2(P)=0\), consider instead the morphism
from the divisor enhancement to the derived based Picard scheme.
The divisor exact sequence shows that its obstruction map surjects
onto \(H^2(\mathcal O_X)\). Hence its relative tangent has amplitude
\([0,1]\), and the transitivity construction again gives a compatible
virtual-pullback triple.
The absolute virtual dimension of that Picard theory is
\[
 \dim S+1-\chi<0.
\]
Its virtual class is zero, so the restricted divisor virtual class
is zero as well.
For \(y>1\), the condition \(H^2(P)=0\) is automatic: otherwise
duality makes \(K-P\) effective, contrary to ampleness.
For \(y=1\), a nonzero section of \(K-P\) has zero intersection with
\(K\), and therefore has empty divisor. Thus the exceptional line is
exactly \(K\). Chow localization expresses the divisor virtual class
as a pushforward from the canonical linear system \(|K|\).
For \(y=0\), an effective divisor is empty; its tangent complex is
zero and its class is the base fundamental class.

We have reduced the low divisor support to \(0\) or \(K\), and hence
the high line to \(3K\) or \(2K\). On a neighborhood of either support
the high direction is smooth: Kodaira vanishing gives its projective
bundle of sections, together with its smooth point Hilbert scheme.
The relative virtual tangent class there is
\begin{equation}
\label{eq:high-tangent}
 [\mathbf R\pi_*\mathcal O(\mathcal C_{\mathrm{high}})]
 -1-[F_{\mathrm{high}}]+[T_{X^{[r_{\mathrm{high}}]}/S}].
\end{equation}
Indeed, the unrestricted rank-one arrow class subtracts
\(\mathbf R\pi_*\mathcal O_X\), while fixing the based determinant
subtracts \(\mathbf R\pi_*\mathcal I_o[1]\); together these subtract
exactly one. The remaining point-ideal calculation is that in
Lemma~\ref{lem:point-excess}.
The projective Euler sequence identifies actual tangent minus
\eqref{eq:high-tangent} with \([F_{\mathrm{high}}]\).

For a smooth morphism with a relative perfect obstruction theory,
the obstruction sheaf is locally free and its virtual pullback is
smooth pullback capped with its top Chern class. This follows by
locally splitting the cotangent comparison onto the locally free
relative cotangent; equivalently, the intrinsic cone is
\([0/T_{\mathrm{rel}}]\) and zero-section excess gives the Euler
factor. The argument works on arbitrary base cycles, not just the
fundamental cycle; compare \cite[Proposition 7.3]{BF} and
\cite[Theorem 4.4]{Manolache}.
Here the obstruction bundle has K-class \([F_{\mathrm{high}}]\),
so its Euler class is \(e(F_{\mathrm{high}})\).  Let
\(\mathcal D\subset\operatorname{Div}_{P}(X/S)\) be the empty-divisor
locus \(S\) or the canonical linear system \(|K|\).
Base-change \(f\) to
\(\mathcal D\times_S X^{[r_{\mathrm{low}}]}_{/S}\), and use the
same letters \(f,\rho\) for the resulting maps.  For every cycle
\(\alpha\) on this product, the high virtual pullback is therefore
\begin{equation}
\label{eq:high-smooth-pullback}
 f^!(\alpha)=e(F_{\mathrm{high}})\cap f^*(\alpha).
\end{equation}
In particular, apply it to
\(\alpha=e(F_{\mathrm{low}})\cap\rho^*\beta\), where \(\beta\)
is a cycle on \(\mathcal D\) supplied by Chow localization for the
low divisor virtual class.  Bivariant proper-pushforward compatibility
then returns the result to \(\mathcal F\).  All operational
insertions restrict to the actual linear systems on those cycles;
no formula for the low cycle supported on \(|K|\) is needed.
These linear systems are smooth over \(S\): for \(K\), constant
cohomology dimensions follow from the product of curves and base
change, and for \(2K,3K\) from Kodaira vanishing.
Their universal divisor classes are \(D_i+x_i\), where \(x_i\) is
the hyperplane class.

\subsection{Marked localization and the short-length exclusion}

Adjoin \(m=3t_0\) independent surface positions to \(\mathcal Q\).
At each position insert \(\ch_2^{\mathbb G_m}(Q)\), cap with the
smoothly pulled virtual class, and push to \(X^m_{/S}\).
This push is proper. The torus acts trivially on the output.
Before localization the total codimension is \(2m-v\).
An ordinary output component of codimension \(4t_0\) would therefore
have equivariant exponent
\[
 (2m-v)-4t_0=2t_0-v<0.
\]
It is absent. Thus the ordinary codimension-\(4t_0\) part of the
localized fixed sum is zero after setting the equivariant parameter
equal to one.

On a fixed branch the insertion is a sum, for \(i=1,2\), of
\begin{equation}
\label{eq:fixed-insertion}
 [Z_i]+(L_i+w_i)(D_i+x_i)-\frac{(D_i+x_i)^2}{2}.
\end{equation}
Here \([Z_i]\) denotes the incidence cycle on the corresponding
surface position.
Test every surface slot by the even tensor
\begin{equation}
\label{eq:surface-test}
 z_\gamma=\sum_{\alpha,\beta}\gamma_{\alpha\beta}
 w_\alpha\otimes w_\beta.
\end{equation}
This is a tensor test in the coefficient algebra, not a numerical
intersection test alone.
Since a unit curve class pairs to zero with a pure \(W_b\) input,
a line-only insertion that survives this test uses a positive divisor
degree on both curve axes. Hence each such position supplies at least
two ordinary codimensions, from powers of the two curve canonical
classes.

Suppose \(k\) of the positions use line-only insertions.
Extract these \(2k\) divisor codimensions by the projection formula.
The remaining \(m-k\) incidence positions lie in the union of the two
point subschemes, of total length \(r+s\). Their pushed support
therefore has at most \(r+s\) distinct surface positions and has
codimension at least
\[
 2(m-k-r-s)
\]
in those factors. Operational incidence classes vanish away from
their incidence supports, so this support assertion applies to
successive caps with arbitrary low cycle multipliers.
If \(r+s<t_0\), the displayed codimension is strictly larger than
the available \(4t_0-2k\), and the tested contribution is zero.

By \eqref{eq:fixed-lengths}, the four remaining ordered divisor pairs
have lengths
\[
\begin{array}{c|cccc}
 (D_1,D_2)&(0,3K)&(3K,0)&(K,2K)&(2K,K)\\ \hline
 r+s&\nu+3s_K&\nu&\nu&\nu-s_K .
\end{array}
\]
Only \((D_1,D_2)=(0,D)\) survives.
On this branch the low divisor is empty and its virtual class is the
base class. The calculation is on the product of the two point
Hilbert schemes and \(|D|\), with both point Euler factors and the
ordinary fundamental cycle. We have
\begin{equation}
\label{eq:surviving-branch}
 N=r+s=\nu+ds_K,\qquad
 \dim(|D|/S)=\chi+\frac{d(d-1)}2s_K-1=t_0,\qquad
 x_1=0,\quad x_2=x.
\end{equation}

\subsection{A relative marked Hilbert-scheme recursion}

The surviving fixed contribution must be expanded as a Chow class on
\(X^m_{/S}\), with every output retained.  Relative EGL recursion for
surface pushforwards to a base appears in
\cite[Section~2.1, Lemma~2.1 and Proposition~2.2]{Kong2026};
the same work also treats one and two marked curve outputs
\cite[Section~3.3, Lemma~3.4]{Kong2026}.
The construction below retains an arbitrary finite set of surface
outputs.  This is the form that lets the codimension bound identify
the Chow tensor left by localization.

\begin{lemma}[Marked universality]
\label{lem:marked-universality}
Fix finitely many marked output surface factors.
Push forward, over relative point Hilbert schemes and any additional
surface factors, a finite polynomial in Chern characters of universal
point ideals, their derived duals, line classes, relative tangent Chern
classes, and equality diagonals.
The resulting Chow class is a universal finite sum of decorated
equality diagonals among the outputs, multiplied by relative
characteristic pushes from closed surface components.
One fixed polynomial procedure chooses the expression, independently
of the surface family. No marked output is integrated in the procedure.
\end{lemma}

\begin{proof}
For a positive point length \(r\), use the consecutive nested scheme
\[
 \mathcal N_r=X^{[r-1,r]}_{/S}
 \xrightarrow{\ \psi\ }X^{[r]}_{/S},\qquad
 \sigma=(\phi,\rho):\mathcal N_r\longrightarrow
 X^{[r-1]}_{/S}\times_S X .
 \]
It is the quotient projectivization of the smaller universal ideal
\(\mathcal I_{r-1}\). The universal quotient line \(\mathcal L\)
and the residual graph \(j\) give
\begin{equation}
\label{eq:nested-ideal-sequence}
 0\longrightarrow\psi_X^*\mathcal I_r
 \longrightarrow\phi_X^*\mathcal I_{r-1}
 \longrightarrow j_*\mathcal L\longrightarrow0.
\end{equation}
The consecutive nested scheme and point Hilbert schemes are smooth
over \(S\). These are the smooth-surface facts underlying
\cite[Section 1]{EGL}. For the relative assertion, choose disjoint
etale charts about the finite support. Deformations of a finite
subscheme, and of an inclusion of two such subschemes, lift uniquely
under an etale map to the relative plane. The resulting local
deformation functors are products of the plane functors, base-changed
from the ground field to \(S\); their smoothness gives the assertion.
Over the distinct-point locus \(\psi\) has degree \(r\), so
\(\psi_*[\mathcal N_r]=r[X^{[r]}_{/S}]\).
Pullback followed by pushforward and division by \(r\) therefore
preserves the original calculation with all other factors present.

Put \(\zeta=c_1(\mathcal L)\).
In every surface slot, replace the ideal by
\eqref{eq:nested-ideal-sequence} in K-theory.
GRR for the residual graph gives
\[
 \ch(j_*\mathcal L)
 =j_*\bigl(e^\zeta\rho^*\td(T_{X/S})^{-1}\bigr).
\]
The graph is a regular codimension-two embedding with the indicated
normal bundle. Thus its only new ingredients are an equality
diagonal to the residual surface coordinate, \(\zeta\), and relative
tangent classes.
Derived duals give no new ingredients, since
\(\ch_j(E^\vee)=(-1)^j\ch_j(E)\).

Choose a two-term locally free resolution
\[
 0\longrightarrow V_-\longrightarrow V_+
 \longrightarrow\mathcal I_{r-1}\longrightarrow0 .
\]
The resolution exists by the fiberwise projective-dimension bound
and an ample locally free presentation.
The quotient projectivization is the zero locus of
\(V_-^\vee(1)\) on \(\mathbb P(V_+)\).
Its codimension is \(\rk V_-\), by the dimension of the consecutive
nested scheme, so the section is regular.
The projective bundle formula then gives
\begin{equation}
\label{eq:nested-segre}
 \sigma_*(\zeta^k\cap[\mathcal N_r])
 =\left[\frac{c(V_-^\vee)}{c(V_+^\vee)}\right]_k
 \cap[X^{[r-1]}_{/S}\times_S X].
\end{equation}
This is the resolution Segre rule of
\cite[Section 1, Lemma 1.1]{EGL}.
After substituting the universal sequence, the integrand is a
polynomial in \(\zeta\) with coefficients pulled from the smaller
Hilbert scheme, the residual point, and the unchanged factors.
Formula~\eqref{eq:nested-segre} lowers the total Hilbert length.
Repeated application ends at length zero.

Products of surface pushes are first rewritten using separate surface
slots, by flat base change and the projection formula.
The same applies to GRR expressions for point-ideal Hom complexes
and tautological bundles, so it covers all the Euler factors that
will occur here after the finite Chow-degree truncation.
At length zero only equality diagonals and the specified decorations
remain. An equality edge connecting two previously distinct
components identifies their coordinates transversely.
An edge within an already connected component contributes the excess
class \(c_2(T_{X/S})\).
Each connected component meeting outputs gives their equality
diagonal with all its decorations on the common coordinate;
a component meeting no output is integrated relatively over \(X\).
Fixing the order of reductions and equality edges gives one universal
polynomial procedure throughout.
\end{proof}

Apply Lemma~\ref{lem:marked-universality} to the surviving fixed
branch. For the surface product of curves,
all the allowed output decorations expand into powers of the two
curve canonical classes.
Consider a connected equality block containing \(k\) outputs.
For \(k=1\), pure-slot testing requires at least two decorating
codimensions, one on each axis.
For \(k=2\), a bare diagonal on either curve axis pairs two isotropic
inputs and is zero. Decorations are therefore needed on both axes,
in addition to the codimension-two surface diagonal.
For \(k\geq3\), the equality block already costs \(2(k-1)\)
codimensions.
These lower bounds are summarized in Table~\ref{tab:block-costs}.
Every surviving block costs at least \(4/3\) codimensions per output,
with equality only for an undecorated three-output block.

\begin{table}[tb]
\centering
\begin{tabular}{@{}clc@{}}
\toprule
Outputs & Requirement under isotropic testing & Codimension at least\\
\midrule
\(1\) & A divisor on each curve axis & \(2\)\\
\(2\) & Equality diagonal and a divisor on each axis & \(4\)\\
\(k\geq3\) & Surface equality diagonal & \(2(k-1)\)\\
\bottomrule
\end{tabular}
\caption{The codimension budget with all output positions retained.
Only a bare triple attains cost \(4/3\) per output.}
\label{tab:block-costs}
\end{table}

The total selected codimension is \(4t_0\) on \(3t_0\) outputs.
It follows that only undecorated triple blocks occur.
No line-only insertion survives.
There is no remaining positive horizontal codimension: every closed
surface component must have top relative degree, and every base
bundle contributes only its scalar characteristic term.
Integration over \(|D|\) therefore extracts exactly \(x^{t_0}\).
In particular, the scalar multiplying the triple blocks can be
computed over a point with the same numerical characteristic data,
using the fixed polynomial procedure of
Lemma~\ref{lem:marked-universality}.

\subsection{The scalar interface}

To specify this scalar without an independence assumption on diagonal
classes, use the polynomial procedure just fixed.
For a formal surface input \(z\), write a connected component carrying
\(k\) marked inputs and decoration \(\alpha\) as the moment
\[
 M_k(\alpha)=\int_X z^k\alpha.
\]
Give unmarked moments their actual top intersection numbers.
Give marked moments value zero except for \(M_3(1)=1\).
Thus this evaluation retains precisely bare triple blocks.
In the generating series of Section~\ref{sec:coefficient}, the
variable \(q\) records total point length, \(p\) is the exponential
variable for marked positions, and \(x\) is the high hyperplane class.
Let \(a_{N,t_0}\) be the coefficient defined in
\eqref{eq:scalar-coefficient}, with the moments evaluated in this
way and \(N\) as in \eqref{eq:surviving-branch}.
Equivalently, \((3t_0)!\,a_{N,t_0}\) is the sum of the coefficients
of all bare triple partitions in the labeled marked pushforward.
The polynomial procedure is fixed before imposing relations among
the diagonal classes of a particular surface.  Section~\ref{sec:coefficient}
evaluates this same representative: its moment-additivity argument
is what permits numerical surface calculations to determine the
selected scalar.  Thus the passage to those calculations does not
replace the marked Chow identity by numerical universality.

\begin{proposition}[Localization relation]
\label{prop:localization}
Under the hypotheses of Theorem~\ref{thm:surface}, the scalar
\(a_{N,t_0}\) satisfies
\[
 a_{N,t_0}\,T_\gamma^{\,t_0}=0
 \qquad\text{in }\mathfrak B_{S,b-1}.
\]
\end{proposition}

\begin{proof}
The codimension-\(4t_0\) marked localization sum is zero.
The low-divisor and short-length arguments discard every fixed
contribution except the branch \((0,D)\), in the stated support
quotient. Marked universality and the output-degree bound leave only
its bare triple partitions, with total coefficient
\((3t_0)!\,a_{N,t_0}\).

Test each surface position with \(z_\gamma\) from
\eqref{eq:surface-test}. A bare surface triple diagonal is the
product of the two curve triple diagonals, hence evaluates to
\(T_\gamma\) up to the Koszul sign from grouping the two curve axes.
This sign is the same for every block and every triple partition:
each whole surface input is even, so permuting those inputs
introduces no further sign. Thus every partition evaluates to the
same signed \(T_\gamma^{t_0}\).
Dividing the resulting equality by its common sign and by
\((3t_0)!\) proves the assertion.
All equalities were Chow equalities or tensor tests in the faithful
coefficient formalism of Proposition~\ref{prop:tensor-test}.
\end{proof}
